\subsection{Laws of Large Numbers}
% Weak LLN, strong LLN, Borel-Cantelli, Kolmogorov inequalities, sums of independent random variables.

The law of large numbers formalizes the idea that a large sample average is
close to the mean of the distribution.

\begin{perspective} \label{per:sampling-distribution-of-sample-mean} \hlt{Sampling Distribution of a Statistic}\\
If $X_1,\ldots,X_n$ are IID with mean $\mu$ and variance $\sigma^2$, then the statistic $\overline{X}_n\defeq n^{-1}\sum_{i=1}^{n}X_i$ is a random variable whose law $\mathcal{L}(\overline{X}_n)$ is its \hlt{sampling distribution}.\\
By Theorem \ref{thm:sample-mean-and-variance-moments}, $\mathbb{E}\overline{X}_n=\mu$ and $\operatorname{Var}(\overline{X}_n)=\sigma^2/n$, so $\mathcal{L}(\overline{X}_n)$ concentrates around $\mu$ as $n\to\infty$.
\end{perspective}

\begin{theorem} \label{thm:weak-law-large-numbers-finite-variance} \hlt{Weak Law of Large Numbers}\\
Let $(X_i)_{i\geq1}$ be IID real-valued random variables with $\mathbb{E}X_1=\mu$, $\operatorname{Var}(X_1)=\sigma^2<+\infty$. Set $\overline{X}_n\defeq n^{-1}\sum_{i=1}^{n}X_i$.\\
Then $\overline{X}_n\xrightarrow{\mathbb{P}}\mu$, equivalently $\mathbb{P}(\abs{\overline{X}_n-\mu}>\varepsilon)\to0$ for every $\varepsilon>0$.
\begin{proof} For $n\geq1$, $\mathbb{E}\overline{X}_n=n^{-1}\sum_{i=1}^{n}\mathbb{E}X_i=\mu$ and, by independence, $\operatorname{Var}(\overline{X}_n)=n^{-2}\sum_{i=1}^{n}\operatorname{Var}(X_i)=\sigma^2/n$.\\
For every $\varepsilon>0$, Chebyshev gives $\mathbb{P}(\abs{\overline{X}_n-\mu}>\varepsilon)\leq\operatorname{Var}(\overline{X}_n)/\varepsilon^2=\sigma^2/(n\varepsilon^2)\to0$, hence $\overline{X}_n\xrightarrow{\mathbb{P}}\mu$.\\
\hlt{Structural Commentary}: Averaging keeps the mean fixed and divides variance by $n$.\\
\hlt{Key Insight}: Chebyshev turns vanishing variance into convergence in probability.
\end{proof}
\end{theorem}

\begin{example} \label{ex:coin-proportion-weak-law-bound} \hlt{Coin Proportions and a Chebyshev Sample Size}\\
Let $X_i\sim\operatorname{Bernoulli}(p)$ be IID head indicators and $\overline{X}_n\defeq n^{-1}\sum_{i=1}^{n}X_i$.\\
Then $\mathbb{E}\overline{X}_n=p$, $\operatorname{Var}(\overline{X}_n)=p(1-p)/n$, and $\overline{X}_n\xrightarrow{\mathbb{P}}p$. \\
or $p=1/2$, Chebyshev gives $\mathbb{P}(0.4\leq\overline{X}_n\leq0.6)\geq1-1/(4n(0.1)^2)=1-25/n$, at least $0.7$ for $n\geq84$.
\end{example}

\begin{theorem} \label{thm:strong-law-large-numbers-integrable} \hlt{Strong Law of Large Numbers}\\
Let $(X_i)_{i\geq1}$ be IID real-valued random variables with $X_1\in L^1$, $\mu\defeq\mathbb{E}X_1$ $\overline{X}_n\defeq n^{-1}\sum_{i=1}^{n}X_i$.\\
Then $\overline{X}_n\xrightarrow{\mathrm{a.s.}}\mu$, equivalently $\mathbb{P}(\lim_{n\to\infty}\overline{X}_n=\mu)=1$.
\begin{proof} After replacing $X_i$ by $X_i-\mu$, assume $\mu=0$. Set $Y_i\defeq X_i\mathbf{1}_{\{\abs{X_i}\leq i\}}$.\\
Since $\sum_{i=1}^{\infty}\mathbb{P}(X_i\neq Y_i)=\sum_{i=1}^{\infty}\mathbb{P}(\abs{X_1}>i)<+\infty$, Borel--Cantelli gives $X_i=Y_i$ eventually a.s.\\
Hence $n^{-1}\sum_{i=1}^{n}(X_i-Y_i)\to0$ a.s. For $Z_i\defeq Y_i-\mathbb{E}Y_i$, independence gives\\
$\sum_{i=1}^{\infty}\operatorname{Var}(Z_i)/i^2\leq\sum_{i=1}^{\infty}\mathbb{E}(X_1^2\mathbf{1}_{\{\abs{X_1}\leq i\}})/i^2<+\infty$, by Kolmogorov's variance criterion, $n^{-1}\sum_{i=1}^{n}Z_i\to0$ a.s.\\
By dominated convergence, $\mathbb{E}Y_i\to0$, hence Cesaro averaging gives $n^{-1}\sum_{i=1}^{n}\mathbb{E}Y_i\to0$.\\
Combining $X_i-Y_i$, $Z_i$, and $\mathbb{E}Y_i$ gives $\overline{X}_n\to0$ a.s.; adding back $\mu$ completes the proof.\\
\hlt{Structural Commentary}: Truncate tails; centered truncations have summable variance.\\
\hlt{Key Insight}: Replace by an eventually equal truncated sequence, then average the centered part.
\end{proof}
\end{theorem}

\begin{corollary} \label{cor:strong-law-implies-weak-law} \hlt{Strong Implies Weak}\\
Let $(X_i)_{i\geq1}$ be IID real-valued random variables, $X_1\in L^1$, $\mu\defeq\mathbb{E}X_1$, $\overline{X}_n\defeq n^{-1}\sum_{i=1}^{n}X_i$. Then $\overline{X}_n\xrightarrow{\mathbb{P}}\mu$.
\begin{proof} By Theorem \ref{thm:strong-law-large-numbers-integrable}, $\overline{X}_n\to\mu$ a.s. For $\varepsilon>0$, set $Y_n\defeq\mathbf{1}_{\{\abs{\overline{X}_n-\mu}>\varepsilon\}}$.\\
Then $Y_n\to0$ a.s. and $0\leq Y_n\leq1$, so dominated convergence gives $\mathbb{P}(\abs{\overline{X}_n-\mu}>\varepsilon)=\mathbb{E}Y_n\to0$.\\
\hlt{Structural Commentary}: Almost-sure convergence gives pointwise tail indicators tending to zero.\\
\hlt{Key Insight}: Dominate the indicators by $1$ and integrate.
\end{proof}
\end{corollary}

\begin{proposition} \label{prop:sample-variance-consistency} \hlt{Consistency of the Sample Variance}\\
Let $(X_i)_{i\geq1}$ be IID real-valued random variables, $\mathbb{E}X_1=\mu$, $\operatorname{Var}(X_1)=\sigma^2<+\infty$.\\
For $n\geq2$, set $\overline{X}_n\defeq n^{-1}\sum_{i=1}^{n}X_i$ and $S_n^2\defeq(n-1)^{-1}\sum_{i=1}^{n}(X_i-\overline{X}_n)^2$.
Then $S_n^2\xrightarrow{\mathbb{P}}\sigma^2$.
\begin{proof}
Let $m_2\defeq\mathbb{E}X_1^2=\sigma^2+\mu^2<+\infty$. By the strong law, $\overline{X}_n\to\mu$ a.s. and $n^{-1}\sum_{i=1}^{n}X_i^2\to m_2$ a.s.\\
Using $S_n^2=\frac{n}{n-1}(n^{-1}\sum_{i=1}^{n}X_i^2-\overline{X}_n^2)$, continuous mapping gives $S_n^2\to m_2-\mu^2=\sigma^2$ a.s.; hence $S_n^2\xrightarrow{\mathbb{P}}\sigma^2$.\\
\hlt{Structural Commentary}: Sample variance is corrected empirical variance.\\
\hlt{Key Insight}: Apply the strong law to $X_i$ and $X_i^2$, then pass through the variance identity.
\end{proof}
\end{proposition}

\begin{remark} \label{rem:cauchy-averages-and-lln-hypotheses} \hlt{Why Integrability Matters}\\
If $X_1\notin L^1$, the mean $\mathbb{E}X_1$ need not exist and sample averages need not converge to any finite deterministic limit; integrability is the condition that makes a law of large numbers meaningful.
\end{remark}
